\tikzexternaldisable
    \centering
    \tikzsetnextfilename{conj_elemental2_R3}
    \begin{tikzpicture}
    \begin{axis}[
        view={125}{25}, % Rotado para que el sólido "venga" hacia adelante
        width=12cm, height=12cm,
        grid=none,
        axis lines=center,
        xmin=0, xmax=6,
        ymin=0, ymax=6, % Eje de profundidad y proyectado
        zmin=0, zmax=6,
        xtick={1,4}, xticklabels={$a$,$b$},
        ytick=\empty, ztick=\empty,
        xlabel={$x$}, ylabel={$y$}, zlabel={$z$},
        every axis x label/.style={at={(ticklabel* cs:1)},anchor=north east},
        every axis y label/.style={at={(ticklabel* cs:1)},anchor=north west},
        every axis z label/.style={at={(ticklabel* cs:1)},anchor=south},
        clip=false
    ]

    % --- 0. Grilla de fondo (plano xz en y=0) ---
    \addplot3 [surf, domain=0:6, domain y=0:6, samples=6, samples y=6, opacity=0.1, point meta=0, colormap={fondogris}{color=(gray!25) color=(gray!25)}, shader=faceted, faceted color=gray] (x,0,y);

    % --- 1. Dominio D_2 en R2 (Plano xz) ---
    % Relleno de D2
    \addplot3 [surf, domain=1:4, domain y=1.2:3.5, samples=2, opacity=0.15, colormap={gray}{color=(gray!20) color=(gray!20)}, shader=interp, draw=none] 
        (x, 0, y);
    % Contorno de D2 (borde recto)
    \draw[thin, black] (1, 0, 1.2) -- (4, 0, 1.2) -- (4, 0, 3.5) -- (1, 0, 3.5) -- cycle;

    % --- 2. Paredes del volumen (Relleno completo de las 4 caras laterales) ---
    % Pared en z = 3.5 (Toma en cuenta z=3.5 para cerrar perfecto con el techo)
    \addplot3 [surf, domain=1:4, domain y=0:1, samples=20, opacity=0.3, colormap={teal}{color=(teal!40) color=(teal!40)}, shader=interp, draw=black!70]
        (x, { (1-y)*(1.2 + 0.2*sin(x*50) + 0.1*3.5) + y*(4.5 + 0.3*cos(x*40) - 0.2*3.5) }, 3.5);
    % Pared en z = 1.2 (Toma en cuenta z=1.2 para cerrar perfecto con el techo)
    \addplot3 [surf, domain=1:4, domain y=0:1, samples=20, opacity=0.3, colormap={teal}{color=(teal!40) color=(teal!40)}, shader=interp, draw=black!70]
        (x, { (1-y)*(1.2 + 0.2*sin(x*50) + 0.1*1.2) + y*(4.5 + 0.3*cos(x*40) - 0.2*1.2) }, 1.2);
    % Pared en x = a (x=1)
    \addplot3 [surf, domain=1.2:3.5, domain y=0:1, samples=20, opacity=0.3, colormap={teal}{color=(teal!40) color=(teal!40)}, shader=interp, draw=black!70]
        (1, { (1-y)*(1.2 + 0.2*sin(50) + 0.1*x) + y*(4.5 + 0.3*cos(40) - 0.2*x) }, x);
    % Pared en x = b (x=4)
    \addplot3 [surf, domain=1.2:3.5, domain y=0:1, samples=20, opacity=0.3, colormap={teal}{color=(teal!40) color=(teal!40)}, shader=interp, draw=black!70]
        (4, { (1-y)*(1.2 + 0.2*sin(200) + 0.1*x) + y*(4.5 + 0.3*cos(160) - 0.2*x) }, x);

    % --- 3. Superficie Trasera y = eta12 ---
    \addplot3 [surf, domain=1:4, domain y=1.2:3.5, samples=25, opacity=0.6,
               colormap={pureorange}{rgb255=(0,255,140,0) rgb255=(1,255,140,0)}, shader=flat, draw=orange!70!black]
        (x, { 1.2 + 0.2*sin(x*50) + 0.1*y }, y);

    % --- 4. Superficie Delantera y = eta22 (Mirando al frente) ---
    \addplot3 [surf, domain=1:4, domain y=1.2:3.5, samples=25, opacity=0.7,
               colormap={pureblue}{rgb255=(0,30,144,255) rgb255=(1,30,144,255)}, shader=flat, draw=blue!70!black]
        (x, { 4.5 + 0.3*cos(x*40) - 0.2*y }, y);

    % --- 5. Proyecciones punteadas al plano (una por cada esquina de D2) ---
    \draw[dashed, thin] (1,0,1.2) -- (1, {1.2 + 0.2*sin(50) + 0.1*1.2}, 1.2);
    \draw[dashed, thin] (4,0,1.2) -- (4, {1.2 + 0.2*sin(200) + 0.1*1.2}, 1.2);
    \draw[dashed, thin] (1,0,3.5) -- (1, {1.2 + 0.2*sin(50) + 0.1*3.5}, 3.5);
    \draw[dashed, thin] (4,0,3.5) -- (4, {1.2 + 0.2*sin(200) + 0.1*3.5}, 3.5);

    % --- 6. Etiquetas y flechas (ajustadas para que no tapen la figura) ---
    \node at (2.5, 0, 2.35) {$D_2$};
    
    % Flecha para la cara delantera (techo)
    \node[anchor=south west] at (3.5, 5.5, 4.5) {$y = \eta_{22}(x,z)$};
    \draw[->, thin] (3.5, 5.4, 4.4) to[out=200, in=60] (2.8, 4.1, 2.8);
    
    % Flecha para la cara trasera (piso curvo)
    \node[anchor=north west] at (3.8, 0.5, 0.5) {$y = \eta_{12}(x,z)$};
    \draw[->, thin] (3.8, 0.7, 0.6) to[out=135, in=300] (2.2, 1.7, 1.8);

    \end{axis}
    \end{tikzpicture}
    \caption{Conjunto elemental de Tipo 2.}

